\begin{abstract}
Finite-memory tracking under drift is governed by a coherence-delay trade-off:
retaining more past information reduces statistical noise, but it also pushes
the estimate further away from the present. As the target distribution moves,
each additional sample becomes at once useful and increasingly old, so the
problem is not simply how much to remember, but how long remembered evidence
should remain operationally current.

We study this trade-off under Wasserstein nonstationarity. For the finite-memory
averaging class analyzed here, the tracking error splits into a statistical term
that decreases with memory and a staleness term that grows with drift. Balancing
those two forces yields a \textbf{cube-root law}: the \textbf{optimal memory
horizon} scales as $\boldsymbol{n^* \asymp (C_K/\zeta)^{2/3}}$, while the
corresponding \textbf{minimum tracking error} scales as
$\boldsymbol{\mathcal{E}_{\min} \asymp C_K^{2/3}\zeta^{1/3}}$. In the
critical-window regime, we also prove a minimax lower bound, showing that this
finite-memory floor is structural rather than estimator-specific.

The experiments follow the same pattern. Gaussian drift experiments produce the
predicted `U`-curve, sliding-window and EWMA estimators recover the cube-root
scaling closely, and the drift-regime benchmarks show that the optimal horizon
changes over time. A cube-root cap then turns the law into an operational
horizon regulator, keeping the effective window near the oracle scale across
drift regimes. A final gap-cost analysis shows that horizon misalignment induces
nonlinear predictive cost, but that the transfer from misalignment to error is
method-dependent and regime-dependent.

The resulting picture is theorem-first and operational: under drift, useful
memory obeys a finite temporal horizon, and the predictive value of that horizon
depends on the interaction between the drift geometry and the memory policy.
Viewed through an Age-of-Information lens, the law quantifies how much freshness
a finite-memory learner can preserve while tracking a changing world.

\end{abstract}

\bigskip
\noindent\textbf{Keywords:}
tracking under drift; finite-memory estimation; Wasserstein distance;
Sinkhorn divergence; concept drift; information freshness;
Age of Information; adaptive memory.

\noindent\textbf{MSC~2020:} 93D05, 93D30, 68P99.
